% =============================================================================
% AGENTE  : Cientista-Líder + Statistician (descrição do pipeline)
% MÓDULO  : Material e Métodos — Português Brasileiro (methods_pt.tex)
% SINCRON.: Deve ser consistente com methods_en.tex
% STATUS  : RASCUNHO
% =============================================================================

\section{Material e Métodos}
\label{sec:methods}

\subsection{Simulação de Dados de Marcadores}
\label{sec:sim}

Para demonstrar o pipeline analítico, um conjunto de dados composto por
60 marcadores moleculares foi simulado utilizando a linguagem de programação
Python (v3.10+) com as bibliotecas \texttt{pandas} (v2.x) e
\texttt{numpy} (v1.24+). A simulação compreendeu 36 marcadores \snp{} e
24 marcadores \ssr{}.

As frequências do alelo principal ($p$) para os marcadores \snp{} foram
amostradas a partir de uma distribuição Beta com parâmetros de forma
$\alpha = 6$ e $\beta = 2$, gerando frequências alélicas enviesadas em
direção às variantes mais comuns, de forma consistente com observações
empíricas em genomas de plantas cultivadas \autocite{Elshire2011}. Para
os marcadores \ssr{}, as frequências foram obtidas de uma distribuição
Beta simétrica ($\alpha = \beta = 3$), refletindo o maior polimorfismo
esperado para locos multialélicos modelados como equivalentes bialélicos.

\subsection{Parâmetros Estatísticos}
\label{sec:stats}

Para cada marcador, foram calculadas as seguintes estatísticas:

\paragraph{Conteúdo de Informação sobre Polimorfismo.}
Seguindo \textcite{Botstein1980}, o \pic{} foi calculado como:
\begin{equation}
    \text{PIC} = 1 - \left(p^2 + q^2\right) - 2p^2q^2
    \label{eq:pic}
\end{equation}
onde $p$ e $q = 1 - p$ são as frequências dos dois alelos.

\paragraph{Heterozigose Esperada.}
Sob o equilíbrio de Hardy--Weinberg:
\begin{equation}
    H_e = 2pq
    \label{eq:he}
\end{equation}

\paragraph{Heterozigose Observada.}
Simulada pela adição de ruído gaussiano $\varepsilon \sim \mathcal{N}(0; 0,04)$
a $H_e$, limitado ao intervalo $[0, 1]$:
\begin{equation}
    H_o = \text{clip}\left(H_e + \varepsilon,\; 0,\; 1\right)
    \label{eq:ho}
\end{equation}

\subsection{Visualização dos Dados}
\label{sec:viz}

Boxplots comparando \pic{} e \he{} entre as classes de marcadores \snp{} e
\ssr{} foram gerados com a biblioteca \texttt{matplotlib} (v3.7+) e exportados
como arquivos PDF para inclusão direta neste documento. Todas as análises são
totalmente reprodutíveis por meio do script \texttt{scripts/plot\_test.py}
(semente aleatória: 42).
